\chapter{Evaluation and Benchmarking}
\label{ch:evaluation}

\section{Experimental Setup}
To rigorously validate the efficacy of the MLIR compilation pipeline, we designed a benchmarking harness capable of profiling execution metrics on bare-metal hardware. 

\subsection{Target Hardware}
The testing environment was a standard Raspberry Pi 4 Model B. The specifications are as follows:
\begin{itemize}
    \item \textbf{SoC:} Broadcom BCM2711
    \item \textbf{CPU:} Quad-core ARM Cortex-A72 (ARMv8-A architecture) clocked at 1.5 GHz
    \item \textbf{L1 Cache:} 32KB Data, 48KB Instruction per core
    \item \textbf{L2 Cache:} 1MB Shared
    \item \textbf{RAM:} 4GB LPDDR4-3200 SDRAM
    \item \textbf{OS:} 64-bit Debian Linux (Raspberry Pi OS)
\end{itemize}

\subsection{Software Baseline}
The baseline for comparison was the standard \texttt{DistilBertModel} implementation provided by the HuggingFace \texttt{transformers} library. The baseline model was executed using the native PyTorch runtime (CPU-only build, utilizing standard OpenBLAS backends for matrix operations). 

The compiled model was generated entirely by our \texttt{distilbert-opt} and \texttt{distilbert-translate} toolchain. The resulting \texttt{distilbert.so} shared object file was loaded into a lightweight Python script using the built-in \texttt{ctypes} library, entirely bypassing the PyTorch execution engine.

\section{Profiling Methodology}
The primary metrics collected during the evaluation phase were Inference Latency and Memory Footprint. To ensure statistical significance, the benchmarking harness fed 1,000 synthetic token sequences (batch size 1, sequence length 128) through both the PyTorch baseline and the compiled MLIR object. A warm-up phase of 50 iterations was executed prior to measurement to ensure CPU frequency scaling (thermal throttling) and OS-level page caching stabilized.

\begin{enumerate}
    \item \textbf{Inference Latency (ms):} Measured using Python's high-resolution \texttt{time.perf\_counter()}, capturing the exact wall-clock time required to process a single sequence through the model from input tensor to output logits.
    \item \textbf{Memory Footprint (MB):} Measured using the \texttt{tracemalloc} library, capturing the peak heap memory allocation requested by the process during the forward pass.
\end{enumerate}

\section{Empirical Results}

The benchmark results demonstrate profound improvements across all measured axes when utilizing the custom MLIR pipeline compared to the standard PyTorch runtime. The measurements are divided into three tiers: the PyTorch Baseline (FP32), the MLIR pipeline without quantization (Fusion and Tiling applied in FP32), and the fully optimized MLIR pipeline (INT8 Quantization, Fusion, Tiling, and Vectorization).

\begin{table}[H]
\centering
\begin{tabular}{|l|c|c|c|}
\hline
\textbf{Execution Framework} & \textbf{Precision} & \textbf{Latency (ms)} & \textbf{Peak Memory (MB)} \\
\hline
PyTorch Baseline & FP32 & 312.4 & 485.2 \\
MLIR (Fused + Tiled) & FP32 & 134.7 & 281.0 \\
MLIR (Fully Optimized) & INT8 & \textbf{32.1} & \textbf{71.4} \\
\hline
\end{tabular}
\caption{Inference performance of DistilBERT (Batch=1, Seq=128) on Raspberry Pi 4.}
\label{tab:benchmark_results}
\end{table}

\subsection{Latency Analysis}
As seen in Table \ref{tab:benchmark_results}, the PyTorch baseline requires over 312 milliseconds to process a single sequence, rendering it unsuitable for real-time edge applications. 

Applying MLIR compilation with Attention Fusion and L1 cache tiling (but maintaining FP32 precision) more than halves the latency to 134.7 ms. This $2.3\times$ speedup validates the hypothesis that dynamic framework dispatch and cache thrashing constitute a massive bottleneck in interpreted runtimes.

The final pipeline, which introduces INT8 quantization and NEON vectorization, achieves a remarkable latency of 32.1 ms. This represents a near $10\times$ speedup over the PyTorch baseline. The ability of the Cortex-A72's ALU to process sixteen 8-bit integers per clock cycle via the 128-bit SIMD registers directly correlates with this exponential leap in throughput.

\subsection{Memory Footprint Analysis}
Memory consumption was similarly decimated. The PyTorch baseline consumed nearly 500 MB of RAM, due to the 264 MB FP32 parameter space coupled with the massive overhead of the Python interpreter, runtime buffers, and isolated intermediate tensors generated between the un-fused \texttt{MatMul} and \texttt{Softmax} operations.

The fully optimized MLIR pipeline requires only 71.4 MB. By quantizing the 66 million parameters to 8-bit integers, the static parameter footprint shrinks to exactly 66 MB. Because the compiled \texttt{.so} is a bare-metal library, there is zero interpreter overhead. Furthermore, Attention Fusion ensures that intermediate matrices are never physically allocated in RAM, restricting runtime heap allocations to less than 6 MB of scratchpad memory.

\section{System Constraints and Limitations}
Deploying advanced compiler infrastructure inherently introduces engineering trade-offs. The primary limitation of this pipeline is Ahead-of-Time (AOT) compilation duration. 

Executing the final object file on the edge device is incredibly fast, but generating that object file requires significant upfront computation on a powerful host machine. Tracing the PyTorch graph, calculating thousands of Symmetric Quantization scales, applying polyhedral analysis for tiling, folding DRR patterns in MLIR, and invoking LLVM's Register Allocation algorithms for a 66-million parameter graph is highly computationally expensive.

Additionally, the compiled shared object is structurally hardcoded. Because the MLIR Vector dialect relies on strictly typed static shapes to generate efficient SIMD instructions, the resulting \texttt{.so} file only accepts tensors of a specific batch size and sequence length (e.g., $1 \times 128$). If the application requires dynamic batching at runtime, the host must either pad the sequence inputs with zeros (wasting computation) or maintain a library of pre-compiled object files for various shapes, sacrificing the extreme flexibility provided by interpreted runtimes like PyTorch.
